\chapter{Concentration}
\label{chap:concentration}

This appendix collects a few standard concentration inequalities used throughout
the monograph.

\section{Scalar concentration}

\begin{lemma}[Hoeffding's inequality]\label{lem:hoeffding}
Suppose $X_1,\dots,X_n$ are independent, identically distributed random variables
with mean $\mu$. Let $\bar X_n := \frac{1}{n}\sum_{i=1}^n X_i$.
If $X_i\in[b_-,b_+]$ almost surely, then for any $\epsilon>0$,
\[
\Pr\!\left(\bar X_n - \mu \ge \epsilon\right)
\le
\exp\!\left(-\frac{2n\epsilon^2}{(b_+-b_-)^2}\right),
\qquad
\Pr\!\left(\mu-\bar X_n \ge \epsilon\right)
\le
\exp\!\left(-\frac{2n\epsilon^2}{(b_+-b_-)^2}\right).
\]
Equivalently,
\[
\Pr\!\left(|\bar X_n-\mu|\ge \epsilon\right)
\le 2\exp\!\left(-\frac{2n\epsilon^2}{(b_+-b_-)^2}\right).
\]
\end{lemma}

In particular, inverting the two-sided bound yields: with probability at least
$1-\delta$,
\[
|\bar X_n-\mu|
\le
(b_+-b_-)\sqrt{\frac{\log(2/\delta)}{2n}}.
\]

\section{Sub-Gaussian random variables}

\begin{definition}[Sub-Gaussian random variable]\label{def:sub-gauss}
A real-valued random variable $X$ is $\sigma$-sub-Gaussian if for all
$\lambda\in\mathbb{R}$,
\[
\E\!\left[\exp\big(\lambda(X-\E X)\big)\right]
\le
\exp\!\left(\frac{\lambda^2\sigma^2}{2}\right).
\]
\end{definition}

A centered Gaussian $X\sim \mathcal{N}(0,\sigma^2)$ is $\sigma$-sub-Gaussian.

\begin{theorem}[Sub-Gaussian tail bound]\label{them:sub_gaussian_property_1}
If $X$ is $\sigma$-sub-Gaussian, then for all $\epsilon>0$,
\[
\Pr\!\left(X-\E X \ge \epsilon\right)
\le
\exp\!\left(-\frac{\epsilon^2}{2\sigma^2}\right),
\qquad
\Pr\!\left(|X-\E X|\ge \epsilon\right)
\le
2\exp\!\left(-\frac{\epsilon^2}{2\sigma^2}\right).
\]
\end{theorem}

\begin{lemma}[Stability under scaling and sums]\label{lemma:sub_gaussian_property_2}
If $X$ is $\sigma$-sub-Gaussian then for any $c\in\mathbb{R}$, $cX$ is
$|c|\sigma$-sub-Gaussian. If $X_1$ and $X_2$ are independent and
$\sigma_1$- and $\sigma_2$-sub-Gaussian respectively, then $X_1+X_2$ is
$\sqrt{\sigma_1^2+\sigma_2^2}$-sub-Gaussian.
\end{lemma}

\section{Martingale concentration}

\begin{theorem}[Hoeffding--Azuma inequality (conditional sub-Gaussian form)]\label{thm:Azuma}
Let $(\mathcal{F}_t)_{t\ge 0}$ be a filtration and $(X_t)_{t=1}^N$ a martingale
difference sequence, i.e.\ $X_t$ is $\mathcal{F}_t$-measurable and
$\E[X_t\mid \mathcal{F}_{t-1}]=0$.
Assume $X_t$ is conditionally $\sigma_t$-sub-Gaussian given $\mathcal{F}_{t-1}$,
meaning that for all $\lambda\in\mathbb{R}$,
\[
\E\!\left[\exp(\lambda X_t)\mid \mathcal{F}_{t-1}\right]
\le
\exp\!\left(\frac{\lambda^2\sigma_t^2}{2}\right)\quad \text{a.s.}
\]
Then for all $\epsilon>0$,
\[
\Pr\!\left(\sum_{t=1}^N X_t \ge \epsilon\right)
\le
\exp\!\left(-\frac{\epsilon^2}{2\sum_{t=1}^N\sigma_t^2}\right),
\qquad
\Pr\!\left(\left|\sum_{t=1}^N X_t\right| \ge \epsilon\right)
\le
2\exp\!\left(-\frac{\epsilon^2}{2\sum_{t=1}^N\sigma_t^2}\right).
\]
\end{theorem}

\begin{lemma}[Bernstein's inequality (i.i.d.)]\label{lem:bernstein_iid}
Let $X_1,\dots,X_n$ be independent random variables with
$\mu:=\E[X_i]$ and $\Var(X_i)\le \sigma^2$ for all $i$.
Assume $|X_i-\mu|\le b$ almost surely.
Then for all $\epsilon>0$,
\[
\Pr\!\left(\bar X_n-\mu \ge \epsilon\right)
\le
\exp\!\left(
-\frac{n\epsilon^2}{2\sigma^2+\frac{2}{3}b\epsilon}
\right),
\qquad
\Pr\!\left(|\bar X_n-\mu| \ge \epsilon\right)
\le
2\exp\!\left(
-\frac{n\epsilon^2}{2\sigma^2+\frac{2}{3}b\epsilon}
\right).
\]
\end{lemma}

In particular, with probability at least $1-\delta$,
\[
|\bar X_n-\mu|
\le
\sqrt{\frac{2\sigma^2\log(2/\delta)}{n}}
\;+\;
\frac{2b\log(2/\delta)}{3n}.
\]

\begin{lemma}[Freedman's inequality (martingale Bernstein)]\label{lem:freedman}
Let $(X_t)_{t=1}^n$ be a martingale difference sequence w.r.t.\ a filtration
$(\mathcal{F}_t)_{t\ge 0}$ and assume $|X_t|\le M$ almost surely.
Define the predictable quadratic variation
\[
V_n := \sum_{t=1}^n \E[X_t^2\mid \mathcal{F}_{t-1}].
\]
Then for all $t>0$,
\[
\Pr\!\left(\sum_{i=1}^n X_i \ge t\right)
\le
\exp\!\left(
-\frac{t^2}{2\big(V_n + Mt/3\big)}
\right).
\]
\end{lemma}

\section{A discrete distribution concentration bound}

The following bound is a standard application of bounded-differences/McDiarmid
(e.g.\ \citep{McD89}).

\begin{proposition}[Concentration for discrete distributions]\label{app:discrete}
Let $Z$ be a discrete random variable supported on $\{1,\dots,d\}$ with
distribution $q\in\Delta([d])$, and let $\widehat q$ be the empirical distribution
from $N$ i.i.d.\ samples. Then with probability at least $1-\delta$,
\[
\|\widehat q-q\|_2
\le
\frac{1}{\sqrt{N}} + \sqrt{\frac{\log(1/\delta)}{2N}}.
\]
Consequently,
\[
\|\widehat q-q\|_1
\le
\sqrt{d}\left(\frac{1}{\sqrt{N}} + \sqrt{\frac{\log(1/\delta)}{2N}}\right).
\]
\end{proposition}

\section{Self-normalized bounds and least squares}

\begin{lemma}[Self-normalized bound for vector-valued martingales; \citep{abbasi2011improved}]
\label{lemma:self_norm}
Let $(\varepsilon_t)_{t\ge 1}$ be a real-valued process adapted to a filtration
$(\mathcal{F}_t)_{t\ge 0}$ such that $\E[\varepsilon_t\mid \mathcal{F}_{t-1}]=0$
and $\varepsilon_t$ is conditionally $\sigma$-sub-Gaussian given $\mathcal{F}_{t-1}$.
Let $(X_t)_{t\ge 1}$ be an $\mathbb{R}^d$-valued adapted process, and let
$V_0\in\mathbb{R}^{d\times d}$ be positive definite. Define
\[
V_t := V_0 + \sum_{i=1}^t X_iX_i^\top.
\]
Then for any $\delta\in(0,1)$, with probability at least $1-\delta$,
simultaneously for all $t\ge 1$,
\[
\left\|\sum_{i=1}^t X_i\varepsilon_i\right\|_{V_t^{-1}}
\le
\sigma\sqrt{2\log\!\left(\frac{\det(V_t)^{1/2}\det(V_0)^{-1/2}}{\delta}\right)}.
\]
Equivalently, squaring both sides,
\[
\left\|\sum_{i=1}^t X_i\varepsilon_i\right\|_{V_t^{-1}}^2
\le
2\sigma^2\log\!\left(\frac{\det(V_t)^{1/2}\det(V_0)^{-1/2}}{\delta}\right).
\]
\end{lemma}

The following OLS bound above can be derived by writing
$\hat\theta-\theta^\star = \Lambda^{-1}\sum_i x_i\epsilon_i$ and observing that
$\|\hat\theta-\theta^\star\|_\Lambda = \|\Lambda^{-1/2}\sum_i x_i\epsilon_i\|_2$
is the Euclidean norm of a $d$-dimensional sub-Gaussian vector (e.g. see also \citep{Hsu2012RandomDA}).

\begin{theorem}[OLS with a fixed design]\label{thm:ols_fixed_design}
Consider a fixed design $\{x_i\}_{i=1}^N\subset\mathbb{R}^d$ and responses
$y_i = \langle \theta^\star, x_i\rangle + \epsilon_i$, where
$\epsilon_1,\dots,\epsilon_N$ are independent, mean-zero, $\sigma$-sub-Gaussian
random variables. Let
\[
\Lambda := \frac{1}{N}\sum_{i=1}^N x_ix_i^\top
\]
and assume $\Lambda$ is invertible. The ordinary least squares estimator is
\[
\hat\theta := \Lambda^{-1}\left(\frac{1}{N}\sum_{i=1}^N x_i y_i\right).
\]
Then for any $\delta\in(0,1)$, with probability at least $1-\delta$,
\[
\|\hat\theta-\theta^\star\|_{\Lambda}
\le
\sigma \cdot \frac{\sqrt{d}+\sqrt{2\log(1/\delta)}}{\sqrt{N}},
\]
and hence
\[
\|\hat\theta-\theta^\star\|_{\Lambda}^2
\le
\sigma^2 \cdot \frac{d + 2\sqrt{2d\log(1/\delta)} + 2\log(1/\delta)}{N}.
\]
\end{theorem}

The following is a standard ``fast-rate'' bound for squared-loss ERM over a
finite function class (e.g.\ \cite{Gyrfi2002ADT}), which we include for
completeness.

\begin{lemma}[Least-squares generalization bound (finite class)]
\label{lem:least_square_gen}
Let $(X_1,Y_1),\dots,(X_n,Y_n)$ be i.i.d.\ with $X_i\sim \nu$ on $\Xcal$ and
$Y_i\in[0,B]$ almost surely for some $B>0$. Let
\[
f^\star(x) := \E[\,Y \mid X=x\,]
\]
and assume $f^\star(x)\in[0,B]$ for all $x$. Let $\Fcal\subseteq\{f:\Xcal\to[0,B]\}$
be finite and define
\[
\epsilon_{\mathrm{approx}}
:= \min_{f\in\Fcal}\ \|f-f^\star\|_{2,\nu}^2.
\]
Let $\hat f\in\argmin_{f\in\Fcal}\sum_{i=1}^n (f(X_i)-Y_i)^2$ be a squared-loss ERM.
Then for any $\delta\in(0,1)$, with probability at least $1-\delta$,
\[
\|\hat f - f^\star\|_{2,\nu}^2
\;\le\;
\frac{20\,B^2\,\ln\!\big(2|\Fcal|/\delta\big)}{n}
\;+\;
3\,\epsilon_{\mathrm{approx}}.
\]
\end{lemma}

\begin{proof}
For $f\in\Fcal$ define
\[
z_i^f := (f(X_i)-Y_i)^2 - (f^\star(X_i)-Y_i)^2.
\]
A direct expansion gives
\[
z_i^f = (f(X_i)-f^\star(X_i))\big(f(X_i)+f^\star(X_i)-2Y_i\big).
\]
Since $f,f^\star,Y_i\in[0,B]$, we have $|z_i^f|\le B^2$ almost surely. Moreover,
using $\E[Y_i\mid X_i]=f^\star(X_i)$,
\[
\E[z_i^f\mid X_i] = (f(X_i)-f^\star(X_i))^2,
\qquad\text{so}\qquad
\E[z_i^f] = \|f-f^\star\|_{2,\nu}^2 =: u_f.
\]
Finally,
\[
(z_i^f)^2 \le (f(X_i)-f^\star(X_i))^2\cdot (2B)^2
\quad\Rightarrow\quad
\Var(z_i^f)\le \E[(z_i^f)^2]\le 4B^2 u_f.
\]

Let $L:=\ln(2|\Fcal|/\delta)$. Applying Bernstein’s inequality to $z_i^f$ and to
$-z_i^f$, and taking a union bound over $f\in\Fcal$, yields that with probability
at least $1-\delta$, simultaneously for all $f\in\Fcal$,
\begin{equation}\label{eq:bern_all_f}
\left|\frac{1}{n}\sum_{i=1}^n z_i^f - u_f\right|
\le
\sqrt{\frac{8B^2 u_f\,L}{n}}
+
\frac{2B^2 L}{3n}.
\end{equation}
Let $\tilde f\in\argmin_{f\in\Fcal}\|f-f^\star\|_{2,\nu}^2$, so $u_{\tilde f}=\epsilon_{\mathrm{approx}}$.
Since $\hat f$ minimizes the empirical squared loss,
\[
\frac{1}{n}\sum_{i=1}^n z_i^{\hat f}\ \le\ \frac{1}{n}\sum_{i=1}^n z_i^{\tilde f}.
\]
Using \eqref{eq:bern_all_f} twice (once for $\hat f$, once for $\tilde f$), we get
\[
u_{\hat f}
\le
u_{\tilde f}
+
\sqrt{\frac{8B^2 u_{\tilde f}L}{n}}
+\frac{4B^2 L}{3n}
+
\sqrt{\frac{8B^2 u_{\hat f}L}{n}}.
\]
Let $\alpha:=\frac{8B^2 L}{n}$ and $\beta:=\frac{2B^2 L}{3n}$, and write $u:=u_{\hat f}$ and
$\eps:=u_{\tilde f}=\epsilon_{\mathrm{approx}}$. Then
\[
u \le \eps + \sqrt{\alpha\eps} + 2\beta + \sqrt{\alpha u}.
\]
Using $\sqrt{\alpha u}\le \frac{u}{2}+\frac{\alpha}{2}$ gives
\[
u \le 2\eps + 2\sqrt{\alpha\eps} + 4\beta + \alpha.
\]
Using $2\sqrt{\alpha\eps}\le \eps+\alpha$ yields
\[
u \le 3\eps + 2\alpha + 4\beta.
\]
Finally,
\[
2\alpha+4\beta
=
\left(16+\frac{8}{3}\right)\frac{B^2 L}{n}
=
\frac{56}{3}\frac{B^2 L}{n}
\le
20\frac{B^2 L}{n},
\]
which completes the proof.
\end{proof}



\section{Other Bounds}

\subsection{Covering numbers and uniformization via $\epsilon$-nets}
\label{app:covering_uniformization}

This section collects a few $\epsilon$-net bounds used in
Chapter~\ref{chap:linear_MDPs}.  We work with the sup norm
$\|f\|_\infty := \sup_{s\in\Scal} |f(s)|$ for functions $f:\Scal\to\R$, the Euclidean norm
$\|\cdot\|_2$ for vectors, and the Frobenius norm $\|\cdot\|_F$ for matrices.

\begin{definition}[$\epsilon$-net and covering number]
Let $(\mathcal{X},d)$ be a metric space and let $\mathcal{F}\subseteq\mathcal{X}$.
A set $\mathcal{N}\subseteq \mathcal{F}$ is an $\epsilon$-net of $\mathcal{F}$ if for every
$f\in\mathcal{F}$ there exists $\tilde f\in\mathcal{N}$ such that $d(f,\tilde f)\le \epsilon$.
The (internal) covering number is
\[
\mathcal{N}(\mathcal{F},\epsilon,d)
:=
\min\Big\{|\mathcal{N}|:\ \mathcal{N}\subseteq \mathcal{F}\ \text{is an $\epsilon$-net of $\mathcal{F}$}\Big\}.
\]
\end{definition}

\begin{lemma}[Volumetric bound for Euclidean balls]\label{lem:volumetric_cover_ball}
Let $\mathcal{B}_2^m(R):=\{x\in\R^m:\ \|x\|_2\le R\}$.
For any $\epsilon\in(0,2R]$,
\[
\mathcal{N}\big(\mathcal{B}_2^m(R),\epsilon,\|\cdot\|_2\big)
\ \le\
\left(1+\frac{2R}{\epsilon}\right)^m.
\]
\end{lemma}

\begin{proof}
Let $\mathcal{S}\subseteq \mathcal{B}_2^m(R)$ be a maximal $\epsilon$-separated set:
$\|x-y\|_2>\epsilon$ for all distinct $x,y\in\mathcal{S}$.
Maximality implies $\mathcal{S}$ is an $\epsilon$-net.

The Euclidean balls $\{x + (\epsilon/2)\mathcal{B}_2^m(1): x\in\mathcal{S}\}$ are disjoint and lie
inside $(R+\epsilon/2)\mathcal{B}_2^m(1)$. Comparing volumes gives
$|\mathcal{S}| \le \big(\frac{R+\epsilon/2}{\epsilon/2}\big)^m = (1+2R/\epsilon)^m$.
\end{proof}

\begin{lemma}[Max and clipping are $\|\cdot\|_\infty$ contractions]\label{lem:max_clip_contraction_app}
Let $\{g_a\}_{a\in\Acal}$ and $\{\tilde g_a\}_{a\in\Acal}$ be real-valued functions on $\Scal$ and
define $G(s):=\max_{a} g_a(s)$ and $\tilde G(s):=\max_{a}\tilde g_a(s)$.
Then
\[
\|G-\tilde G\|_\infty \le \sup_{a\in\Acal}\|g_a-\tilde g_a\|_\infty.
\]
Moreover, for any $M>0$, the clipping map $\mathrm{clip}_{[0,M]}(x):=\min\{M,\max\{0,x\}\}$ satisfies
\[
\|\mathrm{clip}_{[0,M]}\circ u - \mathrm{clip}_{[0,M]}\circ \tilde u\|_\infty \le \|u-\tilde u\|_\infty.
\]
\end{lemma}

\begin{proof}
For any $s$,
\[
G(s)-\tilde G(s)=\max_a g_a(s)-\max_a\tilde g_a(s)
\le \max_a (g_a(s)-\tilde g_a(s))
\le \sup_a \|g_a-\tilde g_a\|_\infty,
\]
and swapping the roles of $g$ and $\tilde g$ yields the two-sided bound.
The clipping bound holds because $\mathrm{clip}_{[0,M]}$ is $1$-Lipschitz on $\R$.
\end{proof}

\begin{lemma}[Covering number for an optimistic value-function class]\label{lem:cover_optimistic_value_class_app}
Assume $\sup_{s,a}\|\phi(s,a)\|_2\le 1$.
Fix parameters $L>0$, $B>0$, $\lambda>0$, and $H\ge 1$.
For $(w,\beta,\Lambda)$ with $\|w\|_2\le L$, $\beta\in[0,B]$, and $\Lambda\succeq \lambda I$,
define $f_{w,\beta,\Lambda}:\Scal\to[0,H]$ by
\[
f_{w,\beta,\Lambda}(s)
:=
\mathrm{clip}_{[0,H]}\!\left(
\max_{a\in\Acal}\left(
w^\top\phi(s,a)\ +\ \beta\sqrt{\phi(s,a)^\top \Lambda^{-1}\phi(s,a)}
\right)
\right).
\]
Let
\[
\Fcal
:=
\{f_{w,\beta,\Lambda}:\ \|w\|_2\le L,\ \beta\in[0,B],\ \Lambda\succeq \lambda I\}.
\]
Then for any $\epsilon\in(0,1]$,
\[
\log \mathcal{N}(\Fcal,\epsilon,\|\cdot\|_\infty)
\ \le\
d\log\!\left(1+\frac{6L}{\epsilon}\right)
+\log\!\left(1+\frac{6B}{\sqrt{\lambda}\,\epsilon}\right)
+\frac{d(d+1)}{2}\log\!\left(1+\frac{18B^2\sqrt d}{\lambda\,\epsilon^2}\right).
\]

\smallskip
\noindent
\emph{Remark (adding a known offset).}
If one replaces $w^\top\phi(s,a)$ inside the max by $u(s,a)+w^\top\phi(s,a)$ for any
fixed bounded function $u:\Scal\times\Acal\to\R$, the same covering bound holds:
$u$ cancels in pairwise differences, and Lemma~\ref{lem:max_clip_contraction_app} still applies.
\end{lemma}

\begin{proof}
Fix two triples $(w,\beta,\Lambda)$ and $(\tilde w,\tilde\beta,\tilde\Lambda)$.
For $(s,a)$ define
\[
g(s,a):=w^\top\phi(s,a)+\beta\sqrt{\phi(s,a)^\top\Lambda^{-1}\phi(s,a)},
\qquad
\tilde g(s,a):=\tilde w^\top\phi(s,a)+\tilde\beta\sqrt{\phi(s,a)^\top\tilde\Lambda^{-1}\phi(s,a)}.
\]
By Lemma~\ref{lem:max_clip_contraction_app},
\[
\|f_{w,\beta,\Lambda}-f_{\tilde w,\tilde\beta,\tilde\Lambda}\|_\infty
\le
\sup_{s,a}|g(s,a)-\tilde g(s,a)|.
\]
For fixed $(s,a)$,
\begin{align*}
|g(s,a)-\tilde g(s,a)|
&\le |(w-\tilde w)^\top\phi(s,a)|
+|\beta-\tilde\beta|\cdot \sqrt{\phi(s,a)^\top\Lambda^{-1}\phi(s,a)} \\
&\quad
+\tilde\beta\cdot\left|\sqrt{\phi(s,a)^\top\Lambda^{-1}\phi(s,a)}
-\sqrt{\phi(s,a)^\top\tilde\Lambda^{-1}\phi(s,a)}\right|.
\end{align*}
The first term is at most $\|w-\tilde w\|_2$ since $\|\phi(s,a)\|_2\le 1$.
For the second term, $\Lambda\succeq \lambda I$ implies
$\phi^\top\Lambda^{-1}\phi \le 1/\lambda$, so it is at most $|\beta-\tilde\beta|/\sqrt{\lambda}$.
For the third term, use $|\sqrt{u}-\sqrt{v}|\le \sqrt{|u-v|}$ for $u,v\ge 0$ and
\[
\left|\phi^\top(\Lambda^{-1}-\tilde\Lambda^{-1})\phi\right|
\le \|\Lambda^{-1}-\tilde\Lambda^{-1}\|_2\|\phi\|_2^2
\le \|\Lambda^{-1}-\tilde\Lambda^{-1}\|_F.
\]
Thus
\[
\sup_{s,a}|g(s,a)-\tilde g(s,a)|
\le
\|w-\tilde w\|_2\ +\ \frac{|\beta-\tilde\beta|}{\sqrt{\lambda}}
\ +\ B\sqrt{\|\Lambda^{-1}-\tilde\Lambda^{-1}\|_F}.
\]

It suffices to build:
(i) an $(\epsilon/3)$-net for $\{w:\|w\|_2\le L\}$ in $\|\cdot\|_2$;
(ii) a $(\sqrt{\lambda}\epsilon/3)$-net for $[0,B]$ in absolute value; and
(iii) a $(\epsilon/(3B))^2$-net for $\{\Lambda^{-1}:\ \Lambda\succeq \lambda I\}$ in $\|\cdot\|_F$.

For (i), Lemma~\ref{lem:volumetric_cover_ball} with $m=d$ gives size at most $(1+6L/\epsilon)^d$.
For (ii), an interval cover gives size at most $1+6B/(\sqrt{\lambda}\epsilon)$.
For (iii), note that $\Lambda^{-1}$ is symmetric and $\Lambda\succeq \lambda I$ implies
$\|\Lambda^{-1}\|_F\le \sqrt d/\lambda$. The set lies in a Frobenius ball of radius
$\sqrt d/\lambda$ inside the $\frac{d(d+1)}{2}$-dimensional vector space of symmetric matrices.
Applying Lemma~\ref{lem:volumetric_cover_ball} with
$m=\frac{d(d+1)}{2}$ and $\epsilon_\Lambda=(\epsilon/(3B))^2$ yields a net of size at most
\[
\left(1+\frac{2(\sqrt d/\lambda)}{(\epsilon/(3B))^2}\right)^{\frac{d(d+1)}{2}}
=
\left(1+\frac{18B^2\sqrt d}{\lambda\,\epsilon^2}\right)^{\frac{d(d+1)}{2}}.
\]
Taking the product net gives an $\epsilon$-net of $\Fcal$.
\end{proof}

\begin{lemma}[Uniform self-normalized bound via an $\epsilon$-net]\label{lem:uniform_self_norm_cover_app}
Let $(\mathcal{F}_t)_{t\ge 0}$ be a filtration and let $(x_t)_{t\ge 0}$ be an adapted sequence
in $\R^d$ with $\|x_t\|_2\le 1$.  Fix $\lambda>0$ and define
\[
\Lambda_n := \lambda I+\sum_{t=0}^{n-1} x_tx_t^\top.
\]

Let $\mathcal{F}$ be a class of functions $f:\Scal\to[0,H]$, and for each $f\in\mathcal{F}$ let
$(\varepsilon_t(f))_{t\ge 0}$ be a sequence such that $\varepsilon_t(f)$ is
$\mathcal{F}_{t+1}$-measurable, $\E[\varepsilon_t(f)\mid \mathcal{F}_t]=0$, and
$|\varepsilon_t(f)|\le H$ almost surely.
Assume moreover that for all $f,g\in\mathcal{F}$,
\begin{equation}\label{eq:eps_lipschitz}
|\varepsilon_t(f)-\varepsilon_t(g)| \le 2\|f-g\|_\infty
\qquad\text{almost surely for all }t.
\end{equation}

Fix $\epsilon\in(0,1]$ and let $\mathcal{N}_\epsilon\subseteq\mathcal{F}$ be an $\epsilon$-net in $\|\cdot\|_\infty$.
Then for any $\delta\in(0,1)$, with probability at least $1-\delta$, simultaneously for all $n\ge 1$
and all $f\in\mathcal{F}$,
\[
\left\|\sum_{t=0}^{n-1} x_t\,\varepsilon_t(f)\right\|_{\Lambda_n^{-1}}
\le
H\sqrt{2\log\!\left(\frac{\det(\Lambda_n)^{1/2}\det(\lambda I)^{-1/2}\,|\mathcal{N}_\epsilon|}{\delta}\right)}
\;+\;
2\epsilon\sqrt{n}.
\]
\end{lemma}

\begin{proof}
Apply Lemma~\ref{lemma:self_norm} (Appendix~\ref{chap:concentration}) to each fixed $\tilde f\in\mathcal{N}_\epsilon$
with $V_0=\lambda I$ and failure probability $\delta/|\mathcal{N}_\epsilon|$.
Since $|\varepsilon_t(\tilde f)|\le H$, each sequence is conditionally $H$-sub-Gaussian.
A union bound implies that with probability at least $1-\delta$, for all $n\ge 1$ and all $\tilde f\in\mathcal{N}_\epsilon$,
\[
\left\|\sum_{t=0}^{n-1} x_t\,\varepsilon_t(\tilde f)\right\|_{\Lambda_n^{-1}}
\le
H\sqrt{2\log\!\left(\frac{\det(\Lambda_n)^{1/2}\det(\lambda I)^{-1/2}\,|\mathcal{N}_\epsilon|}{\delta}\right)}.
\]

Fix any $f\in\mathcal{F}$ and choose $\tilde f\in\mathcal{N}_\epsilon$ with $\|f-\tilde f\|_\infty\le \epsilon$.
Write
\[
\sum_{t=0}^{n-1} x_t\,\varepsilon_t(f)
=
\sum_{t=0}^{n-1} x_t\,\varepsilon_t(\tilde f)
+
\sum_{t=0}^{n-1} x_t\big(\varepsilon_t(f)-\varepsilon_t(\tilde f)\big).
\]
For the second term, define $u_t:=\varepsilon_t(f)-\varepsilon_t(\tilde f)$.
By \eqref{eq:eps_lipschitz}, $|u_t|\le 2\epsilon$, hence $\|u\|_2\le 2\epsilon\sqrt{n}$
for $u=(u_0,\dots,u_{n-1})\in\R^n$.
Let $X=[x_0,\dots,x_{n-1}]\in\R^{d\times n}$ so that $\Lambda_n=\lambda I+XX^\top$ and $\sum_{t=0}^{n-1}x_tu_t=Xu$.
Then
\[
\|Xu\|_{\Lambda_n^{-1}}^2
=
u^\top X^\top (\lambda I+XX^\top)^{-1} X u.
\]
By the Woodbury identity,
$X^\top (\lambda I+XX^\top)^{-1}X = M(I+M)^{-1}$ with $M:=\lambda^{-1}X^\top X\succeq 0$,
so $M(I+M)^{-1}\preceq I$. Therefore $\|Xu\|_{\Lambda_n^{-1}}^2\le \|u\|_2^2$ and
$\|Xu\|_{\Lambda_n^{-1}}\le 2\epsilon\sqrt{n}$.
Combine with the triangle inequality.
\end{proof}

\medskip
We now package the $\epsilon$-net machinery into a single stage-wise uniform confidence statement
tailored to Chapter~\ref{chap:linear_MDPs}.

\begin{corollary}[Uniform confidence over an optimistic value-function class]
\label{cor:uniform_model_conf_bonus_class}
Fix a stage $h$ and an episode budget $K\ge 2$.  Let $(\mathcal{F}_i)_{i\ge 0}$ be the filtration
generated by the interaction up to (and including) $(s_h^i,a_h^i)$ at stage $h$ of episode $i$.
Define
\[
x_i := \phi(s_h^i,a_h^i)\in\R^d,
\qquad
s_{i+1} := s_{h+1}^i\in\Scal,
\qquad
\Lambda_k := \lambda I + \sum_{i=0}^{k-1} x_i x_i^\top,
\]
with $\lambda\ge 1$ and $\|x_i\|_2\le 1$.

Let $\Fcal$ be the optimistic value-function class from Lemma~\ref{lem:cover_optimistic_value_class_app},
so $\Fcal\subseteq\{f:\Scal\to[0,H]\}$.  Assume a linear conditional expectation model at stage $h$:
for every $f\in\Fcal$ there exists $w_f^\star\in\R^d$ such that
\begin{equation}\label{eq:linear_model_stage_h}
\E\!\left[f(s_{i+1})\mid \mathcal{F}_i\right] = x_i^\top w_f^\star
\qquad\text{for all }i\ge 0,
\end{equation}
and $\|w_f^\star\|_2\le H\sqrt d$ for all $f\in\Fcal$.

For each $k\ge 1$ and $f\in\Fcal$, define the ridge estimator
\[
\widehat w_f^k := (\Lambda_k)^{-1}\sum_{i=0}^{k-1} x_i\,f(s_{i+1}),
\qquad
(\widehat P^k f)(s,a) := \phi(s,a)^\top \widehat w_f^k,
\qquad
(P^\star f)(s,a) := \phi(s,a)^\top w_f^\star.
\]

Fix $\delta\in(0,1)$ and set $\epsilon := 1/\sqrt{K}$.
Then with probability at least $1-\delta$, simultaneously for all $k\in\{1,\dots,K\}$,
all $f\in\Fcal$, and all $(s,a)\in\Scal\times\Acal$,
\[
\big|(\widehat P^k f)(s,a) - (P^\star f)(s,a)\big|
\ \le\
\beta_{K,\delta}\;\|\phi(s,a)\|_{(\Lambda_k)^{-1}},
\]
where one may take
\[
\beta_{K,\delta}
:=
H\sqrt{
d\log\!\left(1+\frac{K}{\lambda}\right)
+
2\log\!\left(\frac{\mathcal{N}(\Fcal,\epsilon,\|\cdot\|_\infty)}{\delta}\right)
}
\;+\;
2
\;+\;
\sqrt{\lambda}\,H\sqrt d.
\]
Moreover, by Lemma~\ref{lem:cover_optimistic_value_class_app} and $\epsilon=1/\sqrt K$,
\[
\log \mathcal{N}\!\big(\Fcal,\epsilon,\|\cdot\|_\infty\big)
\le
d\log(1+6L\sqrt K)
+\log\!\Big(1+\frac{6B\sqrt K}{\sqrt{\lambda}}\Big)
+\frac{d(d+1)}{2}\log\!\Big(1+\frac{18B^2\sqrt d\,K}{\lambda}\Big),
\]
so for polynomially bounded $(L,B)$ one has $\beta_{K,\delta}=\widetilde O(Hd)$ as a function
of $(K,1/\delta)$.
\end{corollary}

\begin{proof}
For $f\in\Fcal$, define the noise sequence
\[
\varepsilon_i(f) := f(s_{i+1}) - \E[f(s_{i+1})\mid \mathcal{F}_i].
\]
Then $(\varepsilon_i(f))_{i\ge 0}$ is a martingale difference sequence and
$|\varepsilon_i(f)|\le H$ almost surely.
For any $f,g\in\Fcal$,
\[
|\varepsilon_i(f)-\varepsilon_i(g)|
\le |f(s_{i+1})-g(s_{i+1})|
+ \Big|\E[f(s_{i+1})-g(s_{i+1})\mid \mathcal{F}_i]\Big|
\le 2\|f-g\|_\infty,
\]
so \eqref{eq:eps_lipschitz} holds.

Apply Lemma~\ref{lem:uniform_self_norm_cover_app} to the class $\Fcal$ with this choice of
$\varepsilon_i(\cdot)$ and with $\epsilon=1/\sqrt K$.  On its high-probability event (probability
at least $1-\delta$), for all $k\le K$ and all $f\in\Fcal$,
\[
\left\|\sum_{i=0}^{k-1} x_i\,\varepsilon_i(f)\right\|_{\Lambda_k^{-1}}
\le
H\sqrt{2\log\!\left(\frac{\det(\Lambda_k)^{1/2}\det(\lambda I)^{-1/2}\,\mathcal{N}(\Fcal,\epsilon,\|\cdot\|_\infty)}{\delta}\right)}
+
2\epsilon\sqrt{k}.
\]
Since $\Lambda_k \preceq (\lambda+k)I$, we have
$\det(\Lambda_k)\det(\lambda I)^{-1}\le (1+k/\lambda)^d \le (1+K/\lambda)^d$ and
$2\epsilon\sqrt{k}\le 2$. Hence the above is at most
\[
H\sqrt{d\log\!\left(1+\frac{K}{\lambda}\right)
+
2\log\!\left(\frac{\mathcal{N}(\Fcal,\epsilon,\|\cdot\|_\infty)}{\delta}\right)}
\;+\;2.
\]

Next, using \eqref{eq:linear_model_stage_h}, write $f(s_{i+1}) = x_i^\top w_f^\star + \varepsilon_i(f)$ and compute
\[
\widehat w_f^k - w_f^\star
=
\Lambda_k^{-1}\sum_{i=0}^{k-1} x_i\varepsilon_i(f) \;-\; \lambda\,\Lambda_k^{-1} w_f^\star.
\]
Therefore, for any $(s,a)$,
\begin{align*}
|(\widehat P^k f - P^\star f)(s,a)|
&=
|\phi(s,a)^\top(\widehat w_f^k - w_f^\star)| \\
&\le
\|\phi(s,a)\|_{\Lambda_k^{-1}}\cdot
\left\|\sum_{i=0}^{k-1} x_i\varepsilon_i(f)\right\|_{\Lambda_k^{-1}}
\;+\;
\lambda\,|\phi(s,a)^\top \Lambda_k^{-1}w_f^\star|.
\end{align*}
Since $\Lambda_k\succeq \lambda I$,
\[
\lambda\,|\phi^\top \Lambda_k^{-1}w_f^\star|
\le
\sqrt{\lambda}\,\|\phi\|_{\Lambda_k^{-1}}\,\|w_f^\star\|_2.
\]
Combine the bounds and use $\|w_f^\star\|_2\le H\sqrt d$.
\end{proof}

\subsection{Uniform convergence primitives for Bellman-rank losses}
\label{subsec:br_epsgen_appendix}

This appendix collects standard high-probability uniform convergence bounds that can be used as
$\varepsilon_{\mathrm{gen}}(m,\Hcal,\delta)$ in Assumption~\ref{assum:br_uniform_conv}.
All statements are for a fixed stage $h$ (the chapter union-bounds over $h\in\{0,\dots,H-1\}$).

\paragraph{Notation.}
Fix a loss family $\{\tilde{\ell}(s,a,s';g):g\in\Hcal\}$ and a distribution $\nu$ over $(s,a,s')$.
Given i.i.d.\ data $\Dcal=\{(s_i,a_i,s_i')\}_{i=1}^m\sim \nu^m$, write
\[
\E_{\Dcal}[\tilde{\ell}(\cdot;g)] := \frac{1}{m}\sum_{i=1}^m \tilde{\ell}(s_i,a_i,s_i';g),
\qquad
\E_{\nu}[\tilde{\ell}(\cdot;g)] := \E_{(s,a,s')\sim\nu}[\tilde{\ell}(s,a,s';g)].
\]
We will apply the lemmas below with either
$\tilde{\ell}=\ell_h(\cdot;\,g)$ (the $Q$-version) or
$\tilde{\ell}=\tilde{\ell}_h^{V}(\cdot;\,g)$ (the importance-weighted $V$-version from
\eqref{eq:br_v_loss} in the chapter).

% =========================================================
\subsubsection{Finite hypothesis classes}

\begin{lemma}[Finite class: Hoeffding + union bound]
\label{lem:br_epsgen_finite}
Assume $|\Hcal|<\infty$ and $\|\tilde{\ell}(\cdot;g)\|_\infty\le B$ for all $g\in\Hcal$.
Then with probability at least $1-\delta$,
\[
\sup_{g\in\Hcal}\big|\E_{\Dcal}[\tilde{\ell}(\cdot;g)]-\E_{\nu}[\tilde{\ell}(\cdot;g)]\big|
\ \le\
B\sqrt{\frac{2\ln(2|\Hcal|/\delta)}{m}}.
\]
\end{lemma}

\begin{proof}
Fix $g\in\Hcal$ and apply Hoeffding's inequality to the bounded random variables
$\tilde{\ell}(s_i,a_i,s_i';g)\in[-B,B]$. Union bound over $g\in\Hcal$.
\end{proof}

% =========================================================
\subsubsection{Linear function classes}

We record convenient bounds for Bellman-rank losses induced by linear $Q$-classes.
Fix a feature map $\phi:\Scal\times\Acal\to\R^d$ with $\|\phi(s,a)\|_2\le 1$ for all $(s,a)$.
For parameters $W>0$ and $H\ge 1$, define the stage-wise linear class
\[
\Gcal_h(\phi,W)
:=
\Big\{Q_{h,w}(s,a)=\langle w,\phi(s,a)\rangle:\ \|w\|_2\le W,\ \|Q_{h,w}\|_\infty\le H\Big\}.
\]
At a fixed stage $h$, we will view a hypothesis $g$ as specifying a pair of weights $(w_h,w_{h+1})$,
each lying in the above ball, and use the one-step Bellman residual
\[
\ell_h(s,a,s';g)
:= \langle w_h,\phi(s,a)\rangle - r_h(s,a) - \max_{a'\in\Acal}\langle w_{h+1},\phi(s',a')\rangle.
\]
For the $V$-version with uniform exploration at time $h$, define the importance-weighted loss
\[
\tilde{\ell}_h^{V}(s,a,s';g)
:= \frac{\mathbf{1}\{a=\pi_{h,g}(s)\}}{1/A}\,\ell_h(s,a,s';g),
\qquad
\pi_{h,g}(s)\in\arg\max_{a\in\Acal}\langle w_h,\phi(s,a)\rangle .
\]

\begin{lemma}[Linear classes: uniform convergence for $\ell_h$ and $\tilde{\ell}_h^{V}$]
\label{lem:br_epsgen_linear}
Fix $\delta\in(0,1)$ and let $\Dcal$ be $m$ i.i.d.\ samples from $\nu$.
Assume $r_h(s,a)\in[0,1]$ and $\|Q_{h,w}\|_\infty\le H$ for all $w$ in the class. Then:
\begin{enumerate}
\item (\textbf{$Q$-loss}) With probability at least $1-\delta$,
\[
\sup_{g}\big|\E_{\Dcal}[\ell_h(\cdot;g)]-\E_{\nu}[\ell_h(\cdot;g)]\big|
\ \le\
3H\sqrt{\frac{2d\ln(1+4Wm)+\ln(1/\delta)}{m}}.
\]
\item (\textbf{$V$-loss}) With probability at least $1-\delta$,
\[
\sup_{g}\big|\E_{\Dcal}[\tilde{\ell}_h^{V}(\cdot;g)]-\E_{\nu}[\tilde{\ell}_h^{V}(\cdot;g)]\big|
\ \le\
3AH\sqrt{\frac{2d\ln(1+4AWm)+d\ln(e m A)+\ln(1/\delta)}{m}}.
\]
\end{enumerate}
The suprema are over all stage-local pairs $(w_h,w_{h+1})$ with $\|w_h\|_2,\|w_{h+1}\|_2\le W$ and
$\|Q_{h,w_h}\|_\infty,\|Q_{h+1,w_{h+1}}\|_\infty\le H$.
\end{lemma}

\begin{proof}
We use covering-number arguments.

For the $Q$-loss, note that $|\ell_h|\le |Q_{h,w_h}|+|r_h|+|V_{h+1,w_{h+1}}|\le 2H+1\le 3H$.
Let $\Ncal_\epsilon$ be an $\epsilon$-net of the Euclidean ball $\{w:\|w\|_2\le W\}$; a volumetric bound
gives $|\Ncal_\epsilon|\le (1+W/\epsilon)^d$. Consider the finite set
$\Ncal_\epsilon\times\Ncal_\epsilon$ for the pair $(w_h,w_{h+1})$.
For any fixed net pair $(\bar w_h,\bar w_{h+1})$, Hoeffding's inequality gives
\[
\big|\E_{\Dcal}[\ell_h(\cdot;(\bar w_h,\bar w_{h+1}))]-\E_{\nu}[\ell_h(\cdot;(\bar w_h,\bar w_{h+1}))]\big|
\le
3H\sqrt{\frac{\ln(2/\delta')}{2m}}
\]
with probability at least $1-\delta'$. Union bounding over net pairs and taking $\delta'$ so that the
overall failure probability is $\delta$ yields a bound with
$\ln|\Ncal_\epsilon\times\Ncal_\epsilon|\le 2d\ln(1+W/\epsilon)$.
Finally, if $(w_h,w_{h+1})$ is within $\epsilon$ (in $\|\cdot\|_2$) of $(\bar w_h,\bar w_{h+1})$, then
$\sup_{s,a}|\langle w_h-\bar w_h,\phi(s,a)\rangle|\le \epsilon$ and
$\sup_{s'}|\max_{a'}\langle w_{h+1}-\bar w_{h+1},\phi(s',a')\rangle|\le \epsilon$,
so $|\ell_h(\cdot;g)-\ell_h(\cdot;\bar g)|\le 2\epsilon$ pointwise.
Taking $\epsilon=1/(4m)$ and simplifying constants yields the stated bound.

For the $V$-loss, we again have $|\tilde{\ell}_h^{V}|\le A|\ell_h|\le 3AH$.
The indicator $\mathbf{1}\{a=\pi_{h,g}(s)\}$ depends on $w_h$, so we first control the number of
distinct activation patterns it can induce on the sample.
Fix the sample $\{(s_i,a_i)\}_{i=1}^m$. For each $i$, the event that $a_i$ is greedy under $w_h$ is
\[
\langle w_h,\phi(s_i,a_i)\rangle \ge \langle w_h,\phi(s_i,b)\rangle
\qquad\text{for all } b\in\Acal,
\]
which is described by at most $A-1$ halfspace constraints in $w_h$.
Across $i=1,\dots,m$ there are at most $m(A-1)\le mA$ corresponding hyperplanes, and a standard
arrangement bound implies that $\R^d$ is partitioned into at most $(e m A)^{d}$ regions on which all
indicators $\mathbf{1}\{a_i=\pi_{h,g}(s_i)\}$ are constant. This contributes the additive term
$d\ln(e m A)$.

Condition on any fixed activation pattern of the sample. On this event, the indicator is fixed on
each sample point, and the remaining dependence of $\tilde{\ell}_h^{V}$ on $(w_h,w_{h+1})$ is
$O(A)$-Lipschitz on the sample. We can therefore repeat the same $\epsilon$-net argument as above,
now with an effective net radius $\epsilon\simeq 1/(A m)$, and union bound over both the net pairs and
the (at most) $(e m A)^d$ activation patterns. This yields the stated bound.
\end{proof}

% =========================================================
\subsubsection{Unions of linear classes (feature selection)}

Many representation-selection settings take $\Hcal$ to be a union over candidate feature maps.
Let $\Phi$ be a finite set of feature maps $\phi:\Scal\times\Acal\to\R^s$ with $\|\phi(s,a)\|_2\le 1$.
Define a union hypothesis class $\Hcal=\bigcup_{\phi\in\Phi}\Hcal(\phi)$, where $\Hcal(\phi)$ is the
linear class induced by $\phi$ with $\|w\|_2\le W$ and $\|w^\top\phi\|_\infty\le H$.

\begin{lemma}[Union of linear classes: extra $\log|\Phi|$]
\label{lem:br_epsgen_union_linear}
Under the assumptions of Lemma~\ref{lem:br_epsgen_linear} with dimension $d=s$,
the same bounds hold with $\ln(1/\delta)$ replaced by $\ln(|\Phi|/\delta)$, i.e.\ with probability at
least $1-\delta$,
\[
\sup_{g\in\Hcal}\big|\E_{\Dcal}[\tilde{\ell}(\cdot;g)]-\E_{\nu}[\tilde{\ell}(\cdot;g)]\big|
\ \le\
\varepsilon_{\mathrm{gen}}(m,\Hcal,\delta),
\]
where one may take
\[
\varepsilon_{\mathrm{gen}}(m,\Hcal,\delta)
\ \lesssim\
\begin{cases}
H\sqrt{\dfrac{s\ln(1+Wm)+\ln(|\Phi|/\delta)}{m}} & \text{for the $Q$-loss},\\[10pt]
AH\sqrt{\dfrac{s\ln(1+AWm)+s\ln(e m A)+\ln(|\Phi|/\delta)}{m}} & \text{for the $V$-loss}.
\end{cases}
\]
\end{lemma}

\begin{proof}
Apply Lemma~\ref{lem:br_epsgen_linear} to each fixed $\phi\in\Phi$ with failure probability $\delta/|\Phi|$,
then union bound over $\Phi$.
\end{proof}

\paragraph{Specialization to sparse feature selection.}
If $\Phi$ indexes all $s$-subsets of $[D]$, then $|\Phi|=\binom{D}{s}$ and
$\ln|\Phi|\le s\ln(eD/s)$.
Plugging this into Lemma~\ref{lem:br_epsgen_union_linear} yields dependence polynomial in $s$ and
polylogarithmic in $D$ once substituted into Theorem~\ref{thm:br_main}.